% ==========================================
% BAB I PENDAHULUAN
% ==========================================
\chapter{PENDAHULUAN}
\label{chap:pendahuluan}
% --- Latar Belakang ---
\section{Latar Belakang}
Usaha Mikro, Kecil, dan Menengah (UMKM) merupakan sektor dengan kontribusi terbesar terhadap perekonomian Indonesia. Berdasarkan data \textcite{kemenko2024}, lebih dari 64 juta unit usaha dikategorikan sebagai UMKM dan menyumbang lebih dari 60\% Produk Domestik Bruto (PDB) nasional, serta menyerap sekitar 97\% tenaga kerja. Data ini menunjukkan betapa pentingnya sektor usaha berskala mikro, kecil, maupun menengah dalam menopang stabilitas ekonomi nasional.

Secara yuridis, klasifikasi UMKM diatur dalam Undang-Undang Nomor 20 Tahun 2008 tentang Usaha Mikro, Kecil, dan Menengah. Dalam regulasi tersebut, pengelompokan usaha didasarkan pada kriteria kekayaan bersih atau hasil penjualan tahunan. Usaha Mikro memiliki kekayaan bersih paling banyak Rp50 juta atau penjualan tahunan paling banyak Rp300 juta. Usaha Kecil memiliki kekayaan bersih lebih dari Rp50 juta hingga Rp500 juta atau penjualan tahunan lebih dari Rp300 juta hingga Rp2,5 miliar. Sementara itu, Usaha Menengah memiliki kekayaan bersih lebih dari Rp500 juta hingga Rp10 miliar atau penjualan tahunan lebih dari Rp2,5 miliar hingga Rp50 miliar.

Berdasarkan klasifikasi tersebut, tugas akhir ini akan berfokus pada UKM (Usaha Kecil dan Menengah), yaitu kelompok usaha yang berada pada tingkat perkembangan lebih matang dibanding usaha mikro. UKM umumnya memiliki struktur operasional yang lebih kompleks, kebutuhan pengelolaan keuangan yang lebih besar, serta potensi adopsi teknologi yang lebih tinggi. Kompleksitas ini membuat UKM menghadapi tantangan yang semakin signifikan dalam manajemen keuangan dan pengambilan keputusan, termasuk ketidaktepatan dalam penentuan anggaran proyek, ketiadaan mekanisme validasi anggaran, hingga kesulitan menilai kelayakan ekspansi usaha.

Dalam konteks bisnis modern, perubahan lingkungan usaha yang semakin dinamis menuntut UKM untuk memiliki praktik pengelolaan keuangan yang lebih terstruktur dan adaptif. Untuk mencapai tujuan usaha dalam periode tertentu, UKM perlu menentukan apa yang ingin dicapai serta bagaimana cara mencapainya dengan mempertimbangkan kondisi internal dan kebutuhan pasar. Proses tersebut memerlukan perencanaan finansial yang mencakup \textit{Planning, Budgeting, and Forecasting (PBF)}.

\textit{Planning} merupakan proses penetapan tujuan dan arah strategis usaha, baik dalam jangka pendek maupun jangka menengah. Selanjutnya, \textit{budgeting} dilakukan untuk menerjemahkan rencana tersebut ke dalam bentuk alokasi sumber daya keuangan selama periode tertentu. Anggaran berfungsi sebagai pedoman penggunaan dana serta sebagai alat kontrol dalam pelaksanaan proyek atau kegiatan usaha. Sementara itu, \textit{forecasting} digunakan untuk memperkirakan kinerja keuangan berdasarkan data historis dan kondisi usaha saat ini. Proses ini membantu pemilik atau pengelola UKM dalam memantau pencapaian target, mengidentifikasi potensi penyimpangan, serta mengambil keputusan yang tepat waktu.  Terdapat penelitian yang menunjukkan bahwa perencanaan anggaran yang berbasis prediksi dapat meningkatkan kesiapan organisasi dalam menghadapi perubahan pasar, mengalokasikan sumber daya secara efisien, dan memperkuat ketahanan bisnis \autocite{hariharan2018}.

Namun, dalam praktiknya, banyak UKM belum menerapkan ketiga komponen PBF secara terintegrasi. Perencanaan dan penyusunan anggaran sering kali dilakukan secara sederhana tanpa analisis mendalam terhadap data historis maupun evaluasi atas penggunaan anggaran sebelumnya. Kondisi ini menyebabkan proses perencanaan menjadi kurang berbasis data dan berpotensi menghasilkan alokasi anggaran yang tidak optimal.

Perkembangan teknologi kecerdasan buatan (AI) menghadirkan peluang besar untuk meningkatkan akurasi \textit{forecasting} dan efektivitas perencanaan keuangan. Studi terbaru menunjukkan bahwa teknik berbasis \textit{machine learning} mampu mempelajari pola kompleks, mengolah data dalam skala besar, dan beradaptasi terhadap perubahan pasar jauh dengan lebih baik dibanding metode tradisional \autocite{wasserbacher2022}. Integrasi AI tidak hanya berpotensi meningkatkan akurasi prediksi, tetapi juga mendukung penerapan \textit{Planning, Budgeting, and Forecasting (PBF)} secara lebih terintegrasi. AI dapat membantu proses planning melalui analisis pola kinerja historis, mendukung budgeting dengan penyusunan dan evaluasi alokasi anggaran berbasis data, serta memperkuat forecasting melalui proyeksi yang adaptif terhadap perubahan kondisi usaha.

Meskipun pemanfaatan AI dalam \textit{forecasting} dan perencanaan anggaran berkembang pesat di perusahaan besar, tugas akhir yang secara khusus menganalisis penerapannya pada UKM Indonesia masih sangat terbatas. Berdasarkan kondisi tersebut, tugas akhir ini bertujuan mengembangkan implementasi AI untuk membantu UKM melakukan analisis keuangan, proyeksi, dan penyusunan anggaran secara lebih efektif dan adaptif.
% --- Rumusan Masalah ---
\section{Rumusan Masalah}
Berdasarkan latar belakang yang telah dijelaskan, rumusan masalah dalam tugas akhir ini adalah:
\begin{enumerate}
\item	Bagaimana memanfaatkan data perencanaan dan realisasi anggaran historis sebagai dasar dalam penyusunan anggaran baru?
\item   Bagaimana merancang dan membangun sistem keuangan yang mampu menyusun serta memvalidasi anggaran berdasarkan data historis tersebut?
\item   Bagaimana menerapkan pendekatan AI berbasis Retrieval-Augmented Generation (RAG) untuk meningkatkan kualitas analisis dan penyusunan anggaran dalam sistem yang dikembangkan?
\end{enumerate}

% --- Tujuan ---
\section{Tujuan}
Tujuan yang ingin dicapai dalam pelaksanaan tugas akhir ini adalah sebagai berikut:
\begin{enumerate}
\item   Mengidentifikasi dan menyusun mekanisme pemanfaatan data perencanaan dan realisasi anggaran historis sebagai dasar penyusunan anggaran baru.
\item   Merancang dan mengimplementasikan sistem keuangan yang dapat membantu proses penyusunan serta validasi anggaran secara terstruktur.
\item   Menerapkan pendekatan AI berbasis \textit{Retrieval-Augmented Generation (RAG)} dalam sistem untuk meningkatkan kemampuan analisis dan penyusunan anggaran.
\end{enumerate}

% --- Batasan Masalah ---
\section{Batasan Masalah}
Agar tugas akhir ini terarah dan fokus pada tujuan yang ingin dicapai, maka batasan masalah dalam tugas akhir ini ditetapkan sebagai berikut:
\begin{enumerate}
\item   Tugas akhir ini berfokus pada pengembangan sistem perencanaan anggaran untuk UKM yang mendukung penyusunan dan validasi anggaran berbasis data historis. Lingkup penelitian tidak mencakup fungsi keuangan lain seperti akuntansi lengkap, perpajakan, analisis investasi mendalam, maupun manajemen operasional.
\item   Data yang digunakan dalam pengujian sistem berasal dari data simulasi dan/atau data representatif dari studi kasus UKM (startup binaan Lab SCCIC ITB). Data tidak harus mencerminkan keseluruhan kondisi finansial UKM Indonesia.
\item   Pendekatan kecerdasan buatan dibatasi pada penggunaan \textit{Retrieval-Augmented Generation (RAG)} sebagai mekanisme analisis dan rekomendasi, serta integrasi platform otomasi \textit{low-code} yang digunakan untuk alur dasar seperti pengambilan data, pemrosesan, dan penyajian hasil.
\item   Evaluasi sistem dibatasi pada pengujian fungsionalitas sistem dan relevansi output yang dihasilkan, tanpa mencakup aspek keamanan data tingkat lanjut, kepatuhan hukum, maupun integrasi penuh dengan sistem internal UKM.
\item   Sistem digunakan sebagai alat bantu pengambilan keputusan, bukan sebagai sistem yang secara otomatis menentukan atau mengesahkan anggaran.
\end{enumerate}

% --- Metodologi Pengerjaan TA ---
\section{Metodologi}
Tugas akhir ini menggunakan model proses CRISP-DM(\textit{Cross Industry Standard Process for Data Mining}) sebagai kerangka metodologi pengembangan sistem AI berbasis RAG (\textit{Retrieval-Augmented Generation}) untuk mendukung perencanaan anggaran UKM.
CRISP-DM bersifat \textit{industry-independent} dan terdiri dari enam fase iteratif. Sifat iteratif ini penting karena pengembangan RAG biasanya membutuhkan perulangan perbaikan data, strategi \textit{retrieval}, dan \textit{prompt} berdasarkan hasil evaluasi \textcite{chapman2000}.

\begin{enumerate}

\item \textbf{\textit{Business Understanding}}\\
Tahap ini mencakup pemahaman masalah anggaran UKM, batasan operasional, serta kriteria keberhasilan solusi. Aktivitas yang dilakukan meliputi identifikasi kebutuhan pengguna, perumusan tujuan \textit{data mining/AI}, penentuan \textit{success criteria}, dan penyusunan rencana kerja proyek.
Hasil dari tahapan ini adalah dokumen kebutuhan, pendifinisian target, serta kriteria keberhasilan.

\item \textbf{\textit{Data Understanding}}\\
Tahap ini mencakup pengumpulan data/dokumen, eksplorasi, deskripsi, dan pengecekan kualitas data. Data untuk RAG dapat berupa data historis keuangan usaha, dokumen internal usaha, dan referensi eksternal yang relevan. Hasil dari tahapan ini dapat berupa inventarisasi data, deskripsi struktur data, laporan kualitas data, dan risiko yang mungkin terjadi.

\item \textbf{\textit{Data Preparation}}\\
Tahapan ini meliputi aktivitas menyiapkan data agar siap digunakan oleh model. Aktivitas ini mencakup seleksi data, pembersihan, transformasi, dan pembentukan atribut/meta data. Untuk RAG, \textit{data preparation} spesifik meliputi normalisasi dan \textit{cleaning}, anonimisasi data sensitif, konversi dokumen, \textit{Chunking}, dan pembuatan \textit{embeddings} dan penyimpanan ke \textit{vector database} untuk kebutuhan \textit{retrieval}. Hasil dari tahapan ini adalah dataset yang siap digunakan untuk tahap pemodelan.

\item \textbf{\textit{Modeling}}\\
Tahapan ini mencakup pemilihan teknik, pembangunan \textit{test case}, dan pembangunan model. Model yang dibuat untuk RAG meliputi \textit{retriever dan generator}

\item \textbf{\textit{Evaluation}}\\
Tahap ini mencakup pengecekan hasil terhadap tujuan bisnis, menginterpretasi hasil, dan menentukan tindakan perbaikan. Evaluasi untuk RAG meliputi evaluasi kualitas jawaban, keterlacakan, \textit{retrieval}, dan uji pengguna. Hasil evaluasi digunakan untuk iterasi kembali ke fase \textit{data preparation/modeling}.

\item \textbf{\textit{Deployment}}\\
Tahap ini mencakup implementasi akhir serta rencana penerapan, \textit{monitoring}, dan \textit{maintenance}. Hasil dari tahapan ini adalah sistem RAG yang siap digunakan oleh UKM untuk mendukung perencanaan anggaran.

\end{enumerate}
